\documentclass[12pt,a4paper]{article}

% ── packages ───────────────────────────────────────────────────────────────
\usepackage[utf8]{inputenc}
\usepackage[T1]{fontenc}
\usepackage[english]{babel}
\usepackage{graphicx}
\usepackage{booktabs}
\usepackage{amsmath}
\usepackage{hyperref}
\usepackage{geometry}
\usepackage{caption}
\usepackage{float}

\geometry{margin=2.5cm}

\title{Custom-WikiCS Graph Exploration}
\author{}
\date{\today}

\begin{document}
\maketitle

% ═══════════════════════════════════════════════════════════════════════════
\section{Dataset Overview}
% ═══════════════════════════════════════════════════════════════════════════

The Custom-WikiCS graph is a directed graph derived from Wikipedia computer-science articles.
Each node represents an article, and each directed edge represents a hyperlink from one article to another.
In addition, every node is associated with a 768-dimensional embedding vector generated from its article text.

\begin{itemize}
    \item \textbf{Nodes:} 11\,701
    \item \textbf{Edges:} 291\,039
    \item \textbf{Embedding dimension:} 768
\end{itemize}

% ═══════════════════════════════════════════════════════════════════════════
\section{Degree Distribution}
% ═══════════════════════════════════════════════════════════════════════════

Figure~\ref{fig:degree_dist} shows the degree distributions (total, in-degree, and out-degree)
of the graph.  The distributions are heavily right-skewed: most nodes have low degree while a
few hub pages accumulate thousands of incoming links.

\begin{figure}[H]
    \centering
    \includegraphics[width=\textwidth]{img/degree_distribution.png}
    \caption{Degree distribution histograms (total, in-degree, out-degree).}
    \label{fig:degree_dist}
\end{figure}

\begin{table}[H]
\centering
\caption{Degree statistics.}
\label{tab:degree_stats}
\begin{tabular}{lrrrr}
\toprule
 & \textbf{Min} & \textbf{Max} & \textbf{Mean} & \textbf{Median} \\
\midrule
Total Degree  & 0 & 3\,474 & 49.75 & 14 \\
In-Degree     & 0 & 3\,283 & 24.87 & 4  \\
Out-Degree    & 0 & 316    & 24.87 & 9  \\
\bottomrule
\end{tabular}
\end{table}

\subsection{Log--Log Plot}

Figure~\ref{fig:degree_loglog} presents the degree distribution on log--log axes.
The approximately linear relationship in this scale indicates a power-law--like behaviour,
which is characteristic of scale-free networks.

\begin{figure}[H]
    \centering
    \includegraphics[width=0.65\textwidth]{img/degree_distribution_loglog.png}
    \caption{Degree distribution on log--log scale.}
    \label{fig:degree_loglog}
\end{figure}

% ═══════════════════════════════════════════════════════════════════════════
\section{Degree Assortativity}
% ═══════════════════════════════════════════════════════════════════════════

The \emph{degree assortativity coefficient}~$r$ measures whether high-degree nodes
tend to connect to other high-degree nodes (\emph{assortative}, $r>0$) or to low-degree
nodes (\emph{disassortative}, $r<0$).

\[
    r = -0.0918
\]

The negative coefficient indicates that the graph is \textbf{disassortative}:
high-degree nodes tend to connect to low-degree nodes.
This is consistent with the typical structure of information networks (e.g., Wikipedia),
where a few highly-linked hub articles reference many less popular articles.

% ═══════════════════════════════════════════════════════════════════════════
\section{Topic--Link Structure Assortativity}
% ═══════════════════════════════════════════════════════════════════════════

To assess whether linked articles are also \emph{topically similar},
we compare the embedding vectors of connected node pairs against those of
randomly sampled (unconnected) pairs using three metrics:
Cosine Similarity, Pearson Correlation, and Euclidean Distance.

\subsection{Cosine Similarity}

\begin{figure}[H]
    \centering
    \includegraphics[width=0.8\textwidth]{img/topic_link_cosine.png}
    \caption{Cosine similarity distribution: connected vs.\ random pairs.}
    \label{fig:cosine}
\end{figure}

\subsection{Pearson Correlation}

\begin{figure}[H]
    \centering
    \includegraphics[width=0.8\textwidth]{img/topic_link_pearson.png}
    \caption{Pearson correlation distribution: connected vs.\ random pairs.}
    \label{fig:pearson}
\end{figure}

\subsection{Euclidean Distance}

\begin{figure}[H]
    \centering
    \includegraphics[width=0.8\textwidth]{img/topic_link_euclidean.png}
    \caption{Euclidean distance distribution: connected vs.\ random pairs.}
    \label{fig:euclidean}
\end{figure}

\subsection{Summary}

Table~\ref{tab:topic_link} summarises the three metrics.
Connected pairs consistently show \emph{higher} cosine similarity and Pearson correlation
and \emph{lower} Euclidean distance compared with random pairs.
This confirms that the link structure of the graph is positively correlated with topic similarity.

\begin{table}[H]
\centering
\caption{Topic--link assortativity summary.}
\label{tab:topic_link}
\begin{tabular}{lccccc}
\toprule
\textbf{Metric} & \textbf{Conn.\ Mean} & \textbf{Conn.\ Std} & \textbf{Rand.\ Mean} & \textbf{Rand.\ Std} & \textbf{Diff} \\
\midrule
Cosine Similarity   & 0.6253 & 0.0791 & 0.4990 & 0.0645 & +0.1263 \\
Pearson Correlation & 0.6252 & 0.0791 & 0.4989 & 0.0645 & +0.1263 \\
Euclidean Distance  & 32.1073 & 3.8649 & 37.3866 & 2.8966 & $-5.2792$ \\
\bottomrule
\end{tabular}
\end{table}

% ═══════════════════════════════════════════════════════════════════════════
\section{Local Transitivity (Clustering Coefficient)}
% ═══════════════════════════════════════════════════════════════════════════

The \emph{local clustering coefficient} of a node~$v$ measures how many of $v$'s
neighbours are also neighbours of each other.  High values indicate that $v$ sits
inside a tightly-knit cluster.

\begin{figure}[H]
    \centering
    \includegraphics[width=\textwidth]{img/local_transitivity.png}
    \caption{Local transitivity distribution: all nodes (left) and nodes with CC $> 0$ (right).}
    \label{fig:transitivity}
\end{figure}

\begin{table}[H]
\centering
\caption{Local transitivity statistics.}
\label{tab:cc_stats}
\begin{tabular}{lr}
\toprule
\textbf{Statistic} & \textbf{Value} \\
\midrule
Mean CC       & 0.4527 \\
Median CC     & 0.4286 \\
Min CC        & 0.0000 \\
Max CC        & 1.0000 \\
Std CC        & 0.2991 \\
Nodes with CC $= 0$ & 1\,463 \\
Nodes with CC $> 0$ & 10\,238 \\
Global Transitivity  & 0.2623 \\
\bottomrule
\end{tabular}
\end{table}

The mean local clustering coefficient is $0.4527$, which is relatively high,
indicating that node neighbourhoods tend to be well-connected.
The global transitivity of the graph is $0.2623$.
About $12.5\%$ of nodes have a clustering coefficient of exactly~0;
these are either isolated nodes or nodes whose neighbours share no edges among
themselves.

\end{document}
